% ============================================================================
% COURS 04 — CINÉMATIQUE — PARTIE A
% Source : 5-Cinematique2022.docx
% ============================================================================

\section{Objectifs et mises en situation}
\label{sec:cine-objectifs}

\begin{TLCineMethod}{Objectifs du cours}
Décrire et caractériser les mouvements de tous les points des solides
d'un système : trajectoire, position, vitesse et accélération, sans
étudier les causes mécaniques qui produisent ces mouvements.
\end{TLCineMethod}

\begin{center}
\small
\begin{tabularx}{\linewidth}{>{\bfseries}p{2cm}X c}
\toprule
Je connais & Connaissances attendues & \\
\midrule
& Les définitions des vecteurs position, vitesse et accélération. & $\square$\\
& Les relations de composition des vitesses, rotations et torseurs cinématiques. & $\square$\\
& Le champ des vecteurs vitesse d'un solide indéformable. & $\square$\\
& La loi de vitesse en trapèze. & $\square$\\
\midrule
Je sais & Nommer un mouvement et décrire une trajectoire. & $\square$\\
& Dériver un vecteur dans un repère mobile. & $\square$\\
& Déterminer une vitesse par dérivation ou composition. & $\square$\\
& Déterminer une accélération par dérivation du vecteur vitesse. & $\square$\\
& Passer d'un profil d'accélération à la vitesse et à la position. & $\square$\\
\bottomrule
\end{tabularx}
\end{center}

\subsection{Robot de chirurgie}

La chirurgie robotique impose de maîtriser les vitesses et les
trajectoires. Une mauvaise préparation ou un maniement hésitant peut
allonger l'intervention. Le concepteur doit donc relier les ordres
envoyés aux actionneurs à la cinématique réellement imposée aux
effecteurs.

\begin{center}
\TLCineImage[.42\linewidth]{images/image5.jpeg}\hfill
\TLCineImage[.48\linewidth]{images/image6.png}
\end{center}

\subsection{Prothèse transtibiale}

Une prothèse doit reproduire les amplitudes, vitesses et accélérations
d'un pied réel. La modélisation cinématique permet de vérifier les
mouvements possibles et les niveaux d'accélération transmis à
l'utilisateur.

\TLCineImage[.48\linewidth]{images/image7.jpeg}

\section{Introduction à la cinématique}
\label{sec:cine-introduction}

La cinématique est la partie de la mécanique qui décrit le mouvement,
indépendamment des efforts qui le produisent. Les solides sont supposés
\textbf{indéformables} : pour deux points $A$ et $B$ d'un même solide,
\[
\forall t,\qquad \|\overrightarrow{AB}\|=\mathrm{constante}.
\]

\subsection{Base, repère et référentiel}

Une base orthonormée directe est notée
$\mathcal B=(\vec x,\vec y,\vec z)$. Un repère d'espace associe cette
base à une origine $O$ :
\[
\mathcal R=(O,\vec x,\vec y,\vec z).
\]
Un référentiel comprend un repère d'espace et un repère de temps. En
SII, le temps est commun à tous les référentiels ; on parle donc souvent
simplement de repère de référence.

Dans la suite, un solide et le repère qui lui est lié sont souvent
confondus dans les notations.

\section{Produit vectoriel et produit mixte}
\label{sec:cine-produit}

\subsection{Produit vectoriel}

Le produit vectoriel $\vec U\wedge\vec V$ est perpendiculaire au plan
$(\vec U,\vec V)$, orienté de sorte que
$(\vec U,\vec V,\vec U\wedge\vec V)$ soit direct, et de norme :
\[
\boxed{\|\vec U\wedge\vec V\|
=\|\vec U\|\,\|\vec V\|\,|\sin(\vec U,\vec V)|.}
\]

Dans une base orthonormée directe :
\[
\vec x\wedge\vec y=\vec z,
\quad
\vec y\wedge\vec z=\vec x,
\quad
\vec z\wedge\vec x=\vec y.
\]
L'inversion de l'ordre change le signe.

\begin{TLCineDef}{Propriétés}
\[
\vec U\wedge\vec V=-\vec V\wedge\vec U,
\]
\[
\vec U\wedge(\vec V+\vec W)
=\vec U\wedge\vec V+\vec U\wedge\vec W,
\]
\[
(\lambda\vec U)\wedge(\mu\vec V)
=\lambda\mu(\vec U\wedge\vec V).
\]
Le produit vectoriel est nul lorsque les deux vecteurs sont colinéaires
ou lorsque l'un d'eux est nul.
\end{TLCineDef}

Dans une base orthonormée :
\[
\vec U\wedge\vec V=
\begin{vmatrix}
\vec x&\vec y&\vec z\\
U_x&U_y&U_z\\
V_x&V_y&V_z
\end{vmatrix}.
\]

\begin{TLCineApp}{Calcul de produit vectoriel}
Avec
$\vec V_1=2\vec x+2\vec y-3\vec z$ et
$\vec V_2=-\vec x+6\vec z$, calculer
$\vec V_1\wedge\vec V_2$ puis vérifier son orthogonalité avec les deux
vecteurs de départ.
\end{TLCineApp}

\subsection{Produit mixte}

Le produit mixte de trois vecteurs est le réel :
\[
(\vec U,\vec V,\vec W)
=(\vec U\wedge\vec V)\cdot\vec W.
\]
Sa valeur absolue représente le volume du parallélépipède construit sur
les trois vecteurs. Il est invariant par permutation circulaire et
change de signe lorsqu'on échange deux opérandes.

\section{Torseur cinématique}
\label{sec:cine-torseur}

Le mouvement du solide 1 par rapport au solide 0 est décrit, au point
$A$, par le torseur cinématique :
\[
\boxed{
\left\{\mathcal V_{1/0}\right\}_A
=\left\{\begin{array}{c}
\vec\Omega_{1/0}\\
\vec V(A\in1/0)
\end{array}\right\}_A.}
\]
La résultante $\vec\Omega_{1/0}$ est indépendante du point d'expression.
Le moment cinématique est la vitesse du point considéré.

\begin{TLCineDef}{Inversion du mouvement}
\[
\boxed{\left\{\mathcal V_{2/1}\right\}
=-\left\{\mathcal V_{1/2}\right\}.}
\]
Le point d'expression et la base doivent être précisés pour exploiter
les composantes.
\end{TLCineDef}

\subsection{Changer le point d'expression}

Pour un torseur de résultante $\vec R$ et de moment $\vec M_A$ :
\[
\boxed{\vec M_B=\vec M_A+\overrightarrow{BA}\wedge\vec R.}
\]
Pour le torseur cinématique :
\[
\boxed{\vec V(B\in1/0)=\vec V(A\in1/0)
+\overrightarrow{BA}\wedge\vec\Omega_{1/0}.}
\]
Cette relation constitue le \textbf{champ des vecteurs vitesse} d'un
solide indéformable.

\begin{TLCineApp}{Hélice d'avion}
Une hélice tourne autour de l'axe $(K,\vec w)$ avec
$\vec\Omega_{2/1}=-31\vec w\ \mathrm{rad\,s^{-1}}$. Le point $P$ est
tel que $\overrightarrow{PK}=l\vec u$.
Déterminer $\vec V(P\in2/1)$ en déplaçant le torseur de $K$ vers $P$.
\end{TLCineApp}

\subsection{Opérations sur les torseurs}

Deux torseurs sont égaux si leurs résultantes sont égales et si leurs
moments, exprimés au même point, sont égaux. Pour additionner des
torseurs, il faut les réduire au même point et exprimer les composantes
dans une base commune.

Le comoment de deux torseurs est le réel :
\[
\left\{\mathcal T_1\right\}\otimes
\left\{\mathcal T_2\right\}
=\vec R_1\cdot\vec M_{2,A}
+\vec R_2\cdot\vec M_{1,A}.
\]
Il est indépendant du point de réduction.
